\chapter{Related Work}
\label{ch:related}

\section{Data-movement characterization methodologies}

The direct ancestor of this thesis is DAMOV~\cite{oliveira2021damov}, which
systematized the question ``which functions of a workload are
data-movement-bound, and would near-data processing help?'' for CPUs. Its
three-step recipe---screen the functions that matter; characterize them with
hardware counters plus an architecture-independent locality analysis;
classify them into bottleneck families \emph{derived from the data by
clustering}---is adopted here wholesale, and its central metric, the
last-to-first miss ratio (LFMR: of the requests that miss the first cache,
the fraction that also miss the last and reach DRAM), is retained in
redefined form. DAMOV's validation discipline is also adopted: frozen
thresholds stress-tested on held-out inputs, an independent clustering
algorithm reproducing the classes, and cache-size sweeps confirming class
behaviour. Chapter~\ref{ch:method} enumerates, component by component, the
GPU analogue of each DAMOV element and the three places where the GPU forced
a redesign rather than a translation. Where DAMOV's third step is a
simulator-based intervention experiment (host vs.\ host-with-prefetcher vs.\
NDP across a core-count sweep in ZSim/Ramulator), this thesis substitutes
two \emph{measured} interventions on silicon---independent core- and
memory-clock scaling, and the reuse-distance cache-capacity curve---and
defers the simulated axis to generated configuration
(\S\ref{sec:met-parity}).

On the GPU side, the counter-metric curation used here follows the
NCU-based roofline methodology of Cao et
al.~\cite{cao2023gpudb}, who demonstrated a disciplined path from Nsight
Compute counters to roofline placement for GPU database kernels; this thesis
reuses their metric taxonomy (extended with stall-reason and occupancy axes)
and their practice of measuring, rather than quoting, the roofline ceilings.
The roofline model itself is due to Williams et
al.~\cite{williams2009roofline}. Reuse-distance (stack-distance) analysis
originates with Mattson et al.~\cite{mattson1970stack}; its use here as a
\emph{measured} hit-rate-versus-capacity curve from binary-instrumentation
traces (\S\ref{sec:loc-cdf}) is the GPU translation of DAMOV's
architecture-independent locality step, with one addition the CPU literature
does not need: filtering by memory space, without which register-spill and
scratchpad traffic corrupt the locality signal
(\S\ref{sec:loc-correction}).

\section{Processing-in-memory and near-data systems}

The PiM literature spans three decades; the modern wave is surveyed by Mutlu
et al.~\cite{mutlu2019pim} and Ghose et al.~\cite{ghose2019pim}.
Architectures with direct bearing on this thesis's verdicts include
Tesseract~\cite{ahn2015tesseract} (near-memory graph processing---the
lineage for scatter-capable PiM), PIM-enabled
instructions~\cite{ahn2015pei} (fine-grained offload of cache-defeating
operations), TOM~\cite{hsieh2016tom} (transparent offload of GPU code blocks
to 3D-stacked memory, selected by a compiler bandwidth-benefit analysis---
methodologically the closest GPU-offload ancestor), and the workload-driven
placement study of Boroumand et al.~\cite{boroumand2018google} for consumer
workloads. Commercial near-bank silicon---UPMEM
DPUs~\cite{gomezluna2021upmem}, Samsung HBM-PIM~\cite{kwon2021hbmpim},
SK~hynix GDDR6-AiM~\cite{lee2022aim}---establishes that the substrates this
thesis assigns work to are realizable, not hypothetical. In-storage
processing is grounded in Smart-SSD query offload~\cite{do2013smartssd} and
Biscuit~\cite{gu2016biscuit}. What none of these provide is the
\emph{workload side} for Physical AI: a measured statement of which SLAM
computations and which SLAM data structures want which substrate. That is
the gap this thesis fills.

\section{SLAM systems and their evaluation}

The SLAM algorithmic literature is mature; Cadena et
al.~\cite{cadena2016slam} survey it. The systems measured against in this
thesis's accuracy validation are cuVSLAM itself, as documented in NVIDIA's
technical report~\cite{cuvslam2025}, with ORB-SLAM2/3 as the standard
open-source references~\cite{murartal2017orbslam2,campos2021orbslam3}.
Benchmarking SLAM \emph{as a computational workload} (rather than for
accuracy) was pioneered by SLAMBench~\cite{nardi2015slambench}, which
profiled dense SLAM across CPU/GPU platforms at whole-pipeline granularity.
This thesis differs in three ways: it measures a \emph{production} sparse
V-SLAM stack rather than a research kernel; it descends to per-kernel and
per-data-structure granularity; and its output is a substrate placement, not
a platform comparison. Trajectory-accuracy methodology follows the TUM
benchmark conventions~\cite{sturm2012tum} and the evaluation tutorial of
Zhang and Scaramuzza~\cite{zhang2018tutorial}, with Umeyama
alignment~\cite{umeyama1991}. The datasets used---KITTI~\cite{geiger2013kitti},
EuRoC~\cite{burri2016euroc}, TUM RGB-D~\cite{sturm2012tum},
TUM-VI~\cite{schubert2018tumvi}, and ICL-NUIM~\cite{handa2014iclnuim}---are
the community's standard ground-truth corpora.

\section{GPU instrumentation}

Binary instrumentation of GPU machine code is provided by
NVBit~\cite{villa2019nvbit}, on which this thesis's address tracing is
built (with two extensions: launch-window/kernel-name gating to bound
trace volume, and an allocation-event sidecar that timestamps allocator
activity in the trace's own launch-id clock). GPU simulation and power
modelling---Accel-Sim~\cite{khairy2020accelsim} and
AccelWattch~\cite{kandiah2021accelwattch}---appear in this thesis only as
the target of generated configuration for the follow-on study
(\S\ref{sec:roadmap}); consistent with a standing project rule, no simulated
number is reported in the results chapters.

\section{Positioning}

In one sentence per axis: relative to DAMOV, this thesis is the GPU
re-derivation plus the extension from bottleneck class to substrate verdict;
relative to GPU profiling practice, it adds machine-independent locality and
named-data-structure attribution on top of counters; relative to the PiM
systems literature, it supplies the missing measured workload analysis for a
deployed Physical-AI perception stack; relative to SLAM benchmarking, it is
the first study at the granularity where architectural decisions are
actually made---the kernel and the allocation.
